\section*{Discussion}

\textbf{What a reviewer panel buys.} On a benchmark with known flaws, every three-call condition found more planted flaws than a single reviewer (90--91\% against 81\% for a single rigour reviewer and 86\% for a single combined-checklist reviewer, that is, about ten and five percentage points more), but the three-lens panel found no more than the same lens---or a reviewer holding all three checklists---sampled three times. The recall gain of the panel therefore comes from making several calls, the same mechanism that underlies self-consistency and sampling-and-voting ensembles~\cite{wang2023selfconsistency,li2024moreagents}, and it agrees with evidence that multi-agent debate rarely beats budget-matched sampling~\cite{smit2024mad,wang2024rethinking}. Distinct lenses did change the output in a useful way: at equal recall the panel produced fewer generic critiques, a larger share of critiques matching real flaws and about 40\% fewer high-severity alarms on clean items than the self-ensembles (11.4 vs 19.4 per item), with either merge setting, and it used fewer output tokens. Its letter grade also separated flawed from clean items best, but only while duplicates remained essentially unmerged: its severity score separated them no better than three combined-checklist samples (AUROC 0.84 vs 0.85), and once duplicates were merged every multi-call grade saturated at F. The lenses overlapped heavily (83\% of planted flaw instances, and 91\% of those the panel detected, were found by at least two lenses), partly because their prompts share items and plausibly also because the models we tested apply a general checklist whatever lens they are given; the benchmark cannot separate the two, but in either case the case for a panel rests mainly on focus, with a smaller saving in output tokens and call time, rather than on coverage. For users, this means that the number of reviewer calls, and for a single call a fuller checklist, are the main levers for recall, although extra calls add less when the single reviewer is already strong, and the choice between a panel and a self-ensemble can be made on noise and output cost.

\textbf{What reviewers miss.} With three calls, reviewers reliably caught statistical, confounding, causal and biological flaws (93--100\% per family in every three-call condition), including the false premise that dot-product attention is symmetric; single reviewers were less dependable on some flaw types, missing a third or more of circular-validation flaws and, with the rigour lens alone, of effect-size flaws. By contrast, reviewers found only about one in five false statements about the specific model or its tokenization, regardless of the number of calls. These are premises a reviewer can only check if it knows, or looks up, the model's architecture and training procedure, and our reviewers ran without tools. The estimate rests on four planted statements, two of which (both about scGPT) were never found, so it shows the weakness without measuring its size precisely. Additional reviewer calls did not fix it. We expect grounding to help, but did not test it: requiring the executor to read architecture and tokenization facts from the model configuration in executed code would turn them into checkable outputs rather than assertions, and giving reviewers access to model documentation would let them look such facts up.

\textbf{Merging, agreement and grades.} Whether paraphrased critiques are recognised as duplicates depended mainly on calibration: word-overlap (Jaccard) similarity at an uncalibrated threshold merged almost nothing, while the same method at a calibrated threshold came within 0.02 F1 of LLM adjudication. Agreement between lenses was informative---groups raised by several lenses were about ten times as likely to correspond to a planted flaw---which supports escalating agreed issues. At the same time, the benchmark exposed a property of LLM review that any grading scheme must accommodate: reviewers found high-severity concerns in every clean analysis. A count-based letter grade can rank analyses, but it cannot certify one as acceptable, and once duplicates are merged and agreements escalated it saturates. We therefore treat the grade as an ordinal summary for monitoring and do not use an absolute grade threshold as a default stopping criterion.

\textbf{When to stop.} The stopping calibration showed why convergence rules built on reviewer signals fail with current models: the panel's assessment improved sharply after the first revision and then stayed on a plateau that no absolute threshold would accept, and the only adaptive signal that ever fired---grade stability---merely detected that plateau. The few planted flaws that survived a revision were flaws the panel had not flagged, so no rule based on the panel's output could have caught them, while revisions after the first introduced more new errors than the planted flaws they removed. The pre-specified rule, applied to the runs of both configurations together, selected a budget of four revisions. The margin is one run of 24, and with gpt-5.6-sol alone three revisions did as well; moreover, the rule counts only planted flaws left unresolved, so it favours completeness over minimal editing: a budget of one revision introduced the fewest new errors at the cost of slightly more unresolved planted flaws. We therefore set a fixed budget of four revisions as the default and regard the length of the budget as a choice that belongs to the user. The executed case study, which used the same budget with LLM-adjudicated merging, did not converge either: its panels graded 13 of the 15 reviews of the executed runs F, including the final review of every run. An advancement gate on unresolved critical critiques is only useful once severity is calibrated, which our results suggest it is not.

\textbf{What the case study shows.} Verified execution did what it was designed to do. Every number the executed runs reported came from executed code, and the few that also appear as literals in the scripts are analysis settings; the agents' computations agreed with an independent analysis of the same data, and the panel found and the executors fixed real analytical errors, including a circular degree control that all three runs introduced independently. In the one control run without execution, the executor declined to report results rather than inventing them. Because that run produced neither numbers nor a plan, the case study does not estimate how much verified execution reduces unsupported claims; it shows only that execution makes every reported number checkable. The text-only loops of the stopping calibration show, moreover, that the absence of execution does not by itself prevent unsupported statements: there, revisions credited the original analyses with procedures they never contained. The loop's weaknesses were equally clear. Panels raised more concerns than an executor could address, repeatedly asked for data that did not exist, and never declared an analysis acceptable; executors answered by expanding their scripts (three of the four execution failures occurred in iterations that added large resampling machinery) and, in two runs, by writing the verdict into the code. Two of three runs reached the reference conclusion for TRRUST, but none reported the positive part of the reference result, and one run designed itself into non-identifiability. These observations argue for three additions that the evaluation suggests but did not test: reviewers should be told what data are available so that infeasible requests are recognised as such; conclusions, not only numbers, should be traced to computed quantities, which would expose literal verdicts; and a human should set the scope of an investigation and accept its conclusions. MI-Workbench runs autonomously in the sense that a loop needs no intervention; our results indicate that its conclusions should not be accepted without one.

\textbf{Limitations.} All model evaluations used one provider (OpenAI models through the Codex CLI). The Claude Code and Gemini adapters are implemented and pass offline tests against representative CLI output, but they were not run against live models, so the results characterise one provider's models and do not show whether reviewers from other providers miss the same flaws or overlap as strongly across lenses; merge quality and stopping behaviour may also differ for other model families. All benchmark conditions used the MI-specific reviewer prompts; a domain-agnostic reviewer prompt was not tested, so the contribution of the domain-specific prompt content to recall is not measured. Likewise, the optional generative and knowledge roles (Table~\ref{table_roles}) were not evaluated. Codex adds its own harness instructions and a user-level instruction file to every prompt. Both were constant across conditions, so the comparisons between conditions do not depend on them; because the instruction file asks the model to surface inconsistencies and to push back rather than agree, it may raise the absolute numbers of critiques, flaws found and false alarms compared with calls without it, which we did not test. The benchmark contains 18 synthetic write-ups (12 flawed, 6 clean) whose flaws are detectable from the text; flaws in real analyses can be subtler, and only planted flaws are scored, so unplanted concerns judged substantive may include real problems. The items were written and checked with Claude Opus 5.5, a model from a different provider than the reviewers and judges, so no evaluated model assessed text written by itself; the authoring model may nevertheless introduce stylistic regularities, including idiosyncratic phrasing of model facts, that make some flaws easier or harder to detect than in human-written analyses. Both judges are models from the same provider as the reviewers, and each is also one of the evaluated reviewer configurations; LLM judges can favour outputs similar to their own~\cite{zheng2023judging,wang2024fair}. Neither judge favoured its own model: the second judge was slightly more lenient than the primary judge for all three reviewer configurations, and least so for its own reviewers. The two judges also agreed closely, but neither observation excludes a bias shared by models of one provider. The merge evaluation scores only critiques matched to planted flaws, so over-merging of unrelated concerns is not measured; at the calibrated thresholds the lexical methods grouped more critiques across lenses than LLM adjudication did, the rule that keeps critiques of one reviewer call apart was not varied, and the optional sentence-embedding method was not evaluated. The stopping calibration uses twelve items, two configurations and a single judge, which is the same model as the executor and panel in one configuration and may therefore favour that configuration's revisions; in the absence of data the executor could only revise text, and because its panels merged critiques at the uncalibrated lexical threshold, the executor received largely unmerged feedback and the grade trajectory is specific to that setting, whereas the fixed-budget outcomes and the persistence of high-severity critiques do not depend on the merge. The case study covers one question, one model, one tissue and three executed runs, with large differences between replicates; its evaluation was drafted with LLM assistance and compares conclusions that the runs reached with estimands different from ours. On macOS, verified execution relies on \texttt{sandbox-exec}, which Apple marks as deprecated, and only the \texttt{docker} backend caps memory; the confinement of the \texttt{docker} backend was checked only at the level of the container command, not against a running container. Coordination runs on a single node with SQLite, which was adequate for the up to 20 concurrent runs we tested but is not designed for large deployments; the real-backend test covered one item, one low-effort configuration and at most eight concurrent panels, so provider rate limits were not probed.
